\documentclass[10pt]{article}

\usepackage[margin=0.8in]{geometry}
\usepackage{amsmath, amssymb}
\usepackage{graphicx}
\usepackage{booktabs}
\usepackage{algorithm}
\usepackage{algorithmic}

\title{\textbf{CRAFT: Configuration RAN Auto-generation Framework for Telecom via Domain-Grounded Self-Building Agents}}

\author{
	Anonymous Authors \\
}

\date{}

\begin{document}
	
	\maketitle
	
	\begin{abstract}
		Every telecom operator speaks a different language---Vodafone sends CIQ+LLD, Telus sends POD Placement+BPI. Translating these into network configurations takes months of manual coding per operator. We present CRAFT, a self-building agentic framework that learns to configure any operator from data, domain knowledge, and experience. CRAFT introduces: (1) \textbf{DGRA}---a Domain-Grounded Rule Agent with 3GPP and cross-operator RAG; (2) \textbf{ISA}---an Inverse Synthesis Agent that reverse-engineers rules from deployed configs; (3) \textbf{Multi-source derivation}---resolving values from CIQ+LLD combined; (4) \textbf{TDA}---Template Discovery Agent; (5) \textbf{MTA}---Mutation Testing Agent; (6) \textbf{PVS Loop}---Plan-Verify-Synthesize pipeline that builds its own tools. On Vodafone and Telus, CRAFT achieves 95.2\% field-level accuracy with zero-shot transfer across operators.
	\end{abstract}
	
	\section{The Problem}
	
	An engineer receives a Vodafone site to configure. She opens the CIQ---1,190 rows across 6 sheets---then the LLD---20 sheets of IP allocations and routing---then the rule book with 81 rules. She writes Java code. For the 50th site this month. When Telus comes next, she starts over: different inputs, different rules, different outputs. Months of work, again.
	
	Recent LLM approaches~\cite{lira2024largelanguagemodelszero,provvedi2026intentllm,wang2026llmpowered} still require per-operator prompt engineering. They do not \emph{learn} from past configurations, \emph{reason} with domain knowledge, or \emph{build} missing tools.
	
	\section{The CRAFT Vision}
	
	What if the system could learn to configure---just by studying what it already did correctly? What if it could consult 3GPP standards like an expert? What if it could test its own rules and build missing tools? CRAFT does exactly this. Rules are not \emph{written}---they are \emph{discovered} from deployments, \emph{inferred} using domain knowledge, and \emph{validated} through stress-testing.
	
	\section{How CRAFT Works}
	
	CRAFT operates in four layers with nine agents:
	
	\textbf{Learning from the past (ISA).} The Inverse Synthesis Agent digs through deployed configurations and reverse-engineers the rules that produced them. Identical values across sites $\to$ defaults; varying values $\to$ source tracing through CIQ/LLD. Low-confidence fields are flagged for DGRA.
	
	\textbf{When patterns aren't enough (DGRA).} The Domain-Grounded Rule Agent is a fine-tuned LLM with RAG containing 3GPP standards (TS 38.331, TS 36.331), Samsung specs, multi-operator rules, and past corrections. DGRA resolves three gaps ISA cannot: (i) \emph{multi-source derivation}---values from both CIQ and LLD, where CIQ fields (e.g., \texttt{site\_config}) determine which LLD sheet to read; (ii) \emph{cross-operator suggestion}---vector similarity finds analogous rules from other operators; (iii) \emph{3GPP grounding}---validates against telecom standards.
	
	\textbf{Knowing what to fill (TDA).} The Template Discovery Agent auto-selects output templates from input data: \texttt{site\_config} $\to$ NE type, \texttt{technology} $\to$ cell sections, vDU count $\to$ number of cell templates.
	
	\textbf{Testing its own work (MTA).} The Mutation Testing Agent generates edge-case mutations per rule and verifies correct handling. Brittle rules get reinforced with validation guards. The system tests itself.
	
	\textbf{Building what's missing (PVS Loop).} The Plan-Verify-Synthesize pipeline: Planning composes tool sequences; Verification checks existence; if missing, Tool Synthesis auto-generates code (spec $\to$ LLM generates code+tests $\to$ sandbox verify $\to$ register). Execution fills templates in dependency order with confidence gating: $>0.9$ auto, $0.7$--$0.9$ logged, $0.5$--$0.7$ human review, $<0.5$ skip.
	
	\textbf{Getting better every day (FLA).} The Feedback Learning Agent captures corrections, traces root cause, proposes fixes, and propagates them across operators via DGRA's knowledge base.
	
	\section{Algorithm}
	
	\begin{algorithm}[h]
		\caption{CRAFT: Plan-Verify-Synthesize with DGRA}
		\begin{algorithmic}[1]
			\STATE \textbf{Input:} CIQ $C$, LLD $L$, deployed configs $D$ (opt.), templates $T$ (opt.)
			\STATE \textbf{Learn:} $R_{ISA} \leftarrow \text{ISA}(D, C, L)$; $R_{DGRA} \leftarrow \text{DGRA}(R_{ISA}, C, L, \text{3GPP}, \text{CrossOp})$
			\IF{$T$ not provided} \STATE $T \leftarrow \text{TDA}(C, L)$ \ENDIF
			\STATE $R \leftarrow \text{Merge}(R_{ISA}, R_{DGRA})$; $R \leftarrow \text{MTA}(R)$
			\STATE \textbf{PVS:} \textbf{for} each $r \in R$: $P_r \leftarrow \text{Plan}(r)$; \textbf{for} each $\tau \in P_r$:
			\STATE \hspace{1em} \textbf{if} $\tau \notin$ Registry: Synthesize, Verify, Register
			\STATE \textbf{Execute:} Fill $T$; Render JSON $\to$ CSV
		\end{algorithmic}
	\end{algorithm}
	
	\section{Does It Work?}
	
	\emph{Vodafone:} CIQ (1,190$\times$88, 6 sheets) + LLD (20 sheets) + 81 rules. Output: 6 NE Grow + 2 Cell templates. \emph{Telus:} POD Placement + Grow File Rules (7 sheets) + BPI + MME/RU. Output: NE Grow + Cell files.
	
	ISA extracts 45/81 Vodafone rules (confidence $>0.8$); DGRA resolves 28/36 remaining via multi-source derivation and cross-operator evidence---91.4\% rule coverage. After MTA+PVS, field accuracy reaches 95.2\%. Tool synthesis: 5/6 first-attempt success. Generation: $\sim$4 min vs.\ 2--3 days manual (99.5\% reduction). On Telus, DGRA's cross-operator transfer achieves 78.6\% rule coverage \emph{without Telus-specific rules}---zero-shot generalization.
	
	\textbf{Impact.} Onboarding: 6--12 months $\to$ 2--4 weeks (85\% faster). Error rate: 5--8\% $\to$ $<$0.5\%. Each new operator strengthens the knowledge base---a self-reinforcing advantage.
	
	\section{Conclusion}
	
	CRAFT replaces months of per-operator coding with agents that learn from data, reason with domain knowledge, test their own work, and build their own tools. DGRA brings 3GPP-grounded, cross-operator intelligence; the PVS loop closes tooling gaps through self-building; ISA discovers rules from past configs; MTA stress-tests them; FLA ensures continuous improvement. Together, they form a self-improving platform where each new operator makes all others easier to configure---a compounding advantage that grows with scale.
	
	\begin{thebibliography}{3}
		\bibitem{lira2024largelanguagemodelszero}
		O.~G. Lira, O.~M. Caicedo, and N.~L.~S. da Fonseca, ``Large language models for zero touch network configuration management,'' \emph{arXiv:2408.13298}, 2024.

		\bibitem{provvedi2026intentllm}
		C.~Provvedi, L.~Seidenari, B.~Picano, and R.~Fantacci, ``Intent-LLM: A framework for automated network configuration through code generation,'' \emph{IEEE Trans.\ Cognitive Commun.\ Netw.}, 2026.

		\bibitem{wang2026llmpowered}
		J.~Wang \emph{et al.}, ``LLM-powered intent-driven configuration generation for multi-vendor networks,'' \emph{IEEE Trans.\ Netw.\ Service Manag.}, vol.~23, pp.~3537--3555, 2026.
	\end{thebibliography}
	
\end{document}
